\documentclass{article}

\usepackage[submission]{../../../colm/2026/colm2026_conference}
\usepackage{array}
\usepackage{booktabs}
\usepackage{graphicx}
\usepackage{url}
\usepackage{xurl}
\usepackage{hyperref}

\title{Decode-Position Channel Drift at Long-Decode Reasoning: Budget Tuning as a Partial Cross-Model Remedy}
\author{Anonymous authors\\Paper under double-blind review}

\begin{document}
\maketitle

\begin{abstract}
Quamba2 argues that SSM activations exhibit channel-order preservation and
activation persistence, properties that support offline per-channel
quantization for Mamba-family models. We test the corresponding assumption at
reasoning-scale decode horizons by measuring decode-position channel-set
drift: whether channels in the top 1\% by magnitude early in decode remain in
the same protected set later. Across Granite-4.0-H-Small, Nemotron-3-Nano,
DeepSeek-R1-Distill-Qwen-1.5B, and Falcon-H1, strict top-1\% set-leaving is
0.566234756098, 0.533713200380, 0.670572916667, and 0.673611111111,
respectively, at decode horizons up to 20,000 tokens. Layer-type stratification
shows comparable drift at block outputs for SSM/Mamba, attention, and MoE
blocks, so the Quamba2-style persistence assumption fails at the measured
hybrid-model block-output surface rather than only in a single component type.
We then test W4A16 channel-protection policies with frozen controls. Hard
position or scale switches are actively harmful: M2 and M10 are beaten by
random-bin controls. EMA smoothing removes that random-control-beats failure
but is insufficient at top-1\% budget: M11 recovers only 0.048299284138 median
of the BF16-vs-static gap. Cross-tensor activation+K coupling has signal
relative to random controls but not enough recovery at top-1\%. Budget tuning
is the first partially positive result. M11b top-5 recovers 0.449284091125 on
Granite-Small, but with wide CI95 [-1.300900018791, 1.000790626547]; on
Nemotron, top-5 recovers 0.456736183270 with CI95 [0.345887792388,
0.794986745913] but does not beat static top-10, while top-10 recovers
0.814739798903 and beats static top-10 by 0.220356055286 median. ParoQuant
recovers 0.753776848891 on Granite-Small via rotation, showing that the
problem is not W4A16 impossibility but channel-set protection fragility. A
per-position KL analysis finds sublinear rather than superlinear growth, with
AR decays 0.4387--0.5147, weakening compound-error explanations in this
packet. The resulting contribution is a scoped mechanism paper with a partial
budget-tuning remedy, not a settled deployable quantization recipe.
\end{abstract}

\section{Introduction}

Quamba2 argues that SSM activations have channel-order preservation and
activation persistence, properties that motivate offline per-channel
quantization and channel clustering for Mamba1/Mamba2 backbones
\citep{chiang2025quamba2}. We show that the corresponding protected-set
assumption breaks down at reasoning-scale decode horizons in our W4A16
setting: across three hybrid reasoning LLMs (Granite-4.0-H-Small,
Nemotron-3-Nano, Falcon-H1) and one pure-Transformer reference
(DeepSeek-R1-Distill-Qwen-1.5B), 53--67\% of top-1\% magnitude channels at
calibration positions strictly leave the top-1\% set by the final measured
decode position. This is not merely within-set rank shuffling; it is strict
set membership change. The distinction matters for deployment. Static
protected-channel maps, rotations, per-channel formats, and calibration-time
outlier policies are easiest to engineer when the protected set remains stable.
If channels leave the protected set as decode position increases, then a
policy calibrated early in a trace can be stale later in the same trace.

The deployment pressure is not abstract. Artificial Analysis benchmark-cost
data reported by \citet{bort2025reasoningcosts} put OpenAI o1 at
\$2{,}767.05 to evaluate across seven reasoning-heavy benchmarks, compared
with \$108.85 for GPT-4o on the same suite; the report attributes the
roughly 25$\times$ cost gap to both token expansion and higher per-token
prices. Gartner's March 2026 inference-cost forecast predicts that
trillion-parameter LLM inference will become over 90\% cheaper by 2030 than
in 2025, while warning that advanced reasoning demand can still make total
compute scarce \citep{gartner2026inferencecost}. Commercial market forecasts
are not evidence about model internals, but they indicate scale: one
MarketsandMarkets report projects the AI inference market from \$106.15B in
2025 to \$254.98B in 2030, and another projects edge AI hardware from
\$26.14B in 2025 to \$58.90B in 2030
\citep{marketsandmarkets2025aiinference,marketsandmarkets2025edgeaihardware}.
Long-decode quantization failures therefore matter both for datacenter
reasoning workloads and for eventual edge reasoning deployments.

The scope here is deliberately narrow. For large-scale standard-architecture
reasoning models, recent work such as ParoQuant reports substantial W4A16
quality recovery \citep{liang2025paroquant}. The frontier tested in this
draft is different: small reasoning models, hybrid Mamba-2 or parallel-hybrid
architectures, pure PTQ rather than distillation or quantization-aware
training, and long-decode W4A16 channel-level protection.

The current status is a scoped measurement and mechanism paper with one
partial remedy. The empirical result is stronger than the original
Granite/Mamba-specific framing: decode-position channel drift appears in
Granite hybrid models, a Nemotron-3 hybrid, a pure-Transformer DeepSeek
control, and Falcon-H1's parallel hybrid topology. The negative method result
is also sharper: static unions, hard position bins, hard scale bins, top-1\%
EMA smoothing, top-1\% cross-tensor coupling, DecDEC-style reactive selection,
and stable-core-only protection do not produce a reliable W4A16 channel-set
method. The partial positive result is budget tuning. M11b shows that larger
EMA-smoothed protected budgets can recover a meaningful fraction of the
BF16-vs-static gap, but the best budget differs by model and matched static
top-10 protection is already strong on Nemotron. The paper therefore claims
that budget is load-bearing, not that a fixed top-5 EMA recipe is a settled
cross-model method.

The contribution is:
\begin{itemize}
\item quantitative characterization of decode-position channel-set drift at
  reasoning-scale horizons (up to 20,000 decode tokens) across hybrid Mamba-2
  and pure-Transformer reasoning LLMs, extending DecDEC's short-horizon
  Transformer evidence to the long-decode reasoning regime while contradicting
  the channel-persistence assumption that motivates Quamba2-style static
  ordering;
\item cross-architecture replication on Nemotron-3, DeepSeek-R1-Distill, and
  Falcon-H1 under fixed deterministic traces;
\item strict set-membership decomposition showing that many channels leave
  the protected top-1\% set entirely, rather than merely shuffling within it;
\item mechanism evidence that boundary-discontinuous switching is harmful,
  top-1\% smoothing has signal but insufficient magnitude, cross-tensor
  coupling has signal without enough budget, and per-position KL does not show
  superlinear compound-error growth in the measured Granite packet;
\item a partial remedy: M11b budget-tuned EMA protection recovers 0.4493
  median on Granite-Small at top-5 and 0.8147 median on Nemotron at top-10,
  but the Nemotron top-5 arm does not beat matched static top-10, so the method
  is a budget-sensitivity finding rather than a model-invariant recipe;
\item a defensive SOTA context: ParoQuant's algorithmic Granite baseline
  recovers 0.7538 median through rotation, showing that the paper's negative
  results are scoped to channel-set protection rather than to all W4A16 PTQ.
\end{itemize}
The paper does not claim a deployable kernel/runtime system, a universal
quantization recipe, or completed architecture-general validation.

\section{Related Work}

\paragraph{Mamba quantization and dynamic outliers.}
OutlierMigrate is not the first work to observe dynamic outlier behavior in
Mamba-style models. In vision state-space models, QMamba reports that hidden
state sequences exhibit highly dynamic variations and introduces Temporal
Group Quantization for hidden states \citep{li2025qmamba}. OuroMamba makes the
point even more directly for vision Mamba models: it identifies dynamic
activation-outlier variation across time steps as a failure mode for static
post-training quantization and uses lightweight dynamic outlier detection
during inference \citep{ramachandran2025ouromamba}. These results are important
prior art. They establish that dynamic outliers are already known in the
vision-Mamba setting.\footnote{The phrase ``outlier migration'' is overloaded.
In this paper, the operational measurement is decode-position channel drift:
which activation channels are in the top-magnitude set at different decode
positions. SmoothQuant uses migration for activation-to-weight difficulty
transfer \citep{xiao2023smoothquant}. MoBiQuant uses it for token-set shifts
under different bit widths \citep{wang2026mobiquant}; we reserve our
quantitative claims for the channel-across-decode-position phenomenon measured
in this paper.}

Language-Mamba quantization work has emphasized a different operating regime.
Mamba-PTQ shows that recurrent LLMs exhibit activation outlier channels and
frames outlier-aware quantization as a first step \citep{pierro2024mambaptq}.
Quamba proposes a static 8-bit per-tensor recipe that suppresses maximum input
values to selective SSMs and uses a Hadamard transform for output activations
\citep{chiang2024quamba}. Quamba2 is closer to the assumption tested here:
it explicitly builds on channel-order preservation and activation persistence
for Mamba1/Mamba2 quantization \citep{chiang2025quamba2}. Our contribution is
not that dynamic channel sets exist in every setting; it is that the
Quamba2-style persistence assumption does not hold for the measured
reasoning-scale W4A16 decode horizons. MambaQuant uses variance-aligned rotations and
reports heavy-tailed, channel-uneven distributions in Mamba-family models
\citep{xu2025mambaquant}. Quamba-SE softens the treatment of activation
outliers with multiple adaptive scales but still evaluates a per-value
quantizer rather than a decode-time channel migration policy
\citep{chen2026quambase}. OutlierMigrate asks which regime applies to hybrid
LLMs in long reasoning traces: the static or calibration-time protection
regime common in language deployment recipes, or the dynamic regime already
documented in vision Mamba.

\paragraph{Static protection in LLM deployment.}
Several influential LLM quantization systems protect or transform salient
activation structure using calibration-time statistics or fixed transformations.
SmoothQuant migrates activation difficulty into weights using offline
activation statistics \citep{xiao2023smoothquant}; AWQ protects a small set of
activation-salient weight channels \citep{lin2023awq}; QuaRot removes hidden
state outliers with rotations \citep{ashkboos2024quarot}; KVQuant isolates KV
cache outliers per vector and uses per-channel key quantization
\citep{hooper2024kvquant}; and BlockDialect assigns block-wise number formats
from a learned formatbook \citep{jang2025blockdialect}. These systems are not
claimed to be wrong by our current evidence. The implication is narrower: if
the protected high-magnitude set changes substantially during long decode in
hybrid LLMs, static protection maps may need migration-aware updates or
validation on the specific recurrent/linear-attention architecture. Several
of these systems already include adaptive components, such as KVQuant's
per-vector dense-and-sparse handling or BlockDialect's online activation
format selection; our current evidence motivates stress-testing those designs
on hybrid long decode, not declaring them failed.

\paragraph{Dynamic channel and KV-cache policies.}
Dynamic inference-time policies are closer to the methods we test and fail.
DecDEC dynamically selects salient activation channels at each decode step and
fetches residual compensation from CPU memory \citep{park2025decdec}. DecDEC
also reports low recall for static outlier-channel identification at short
decode horizons. Our measurement is therefore not a first demonstration that
channel sets change during decode; the defensible distinction is horizon,
model class, architecture, and benchmark surface: up to 20,000 decode tokens
on reasoning LLMs, including hybrid Mamba-2 and parallel-hybrid models, rather
than DecDEC's short-horizon Transformer setting. M2 and hard-binned M10 differ
from DecDEC in that they switch a precomputed protected-channel or scale policy
by decode-position bin; M11 tests whether smoothing that decision surface
avoids the boundary discontinuities exposed by M2 and M10. It does avoid the
random-control-beats pattern but still does not recover enough quality.
KV-cache systems address a
different tensor. KIVI quantizes key cache per channel and value cache per
token \citep{liu2024kivi}; PM-KVQ targets long-CoT KV-cache quantization and
RoPE-induced key-cache distribution shift \citep{liu2025pmkvq}; and
AttentionPredictor learns to predict future attention scores for KV-cache
compression and prefetch \citep{yang2025attentionpredictor}. Newer
long-context KV-cache systems sharpen the distinction. ChanMix reallocates
bits across sensitive KV-cache channels \citep{liao2026chanmix}, while MixKVQ
uses query-aware mixed precision for critical key channels in long-context
reasoning \citep{zhang2025mixkvq}. These are close neighbors on the
channel-wise long-context quantization axis, but they act on cached key/value
tensors. Our measurements and interventions target linear-layer activation
channels, so a ChanMix- or MixKVQ-style success would not by itself resolve
the W4A16 activation-channel failure measured here. These works motivate
long-decode inference adaptation, but our measurements are about linear-layer
activation channel sets.\footnote{``Hot channels'' in HCP/CHON
(arXiv:2602.02047) refers to persistently extreme channels during NVFP4
pretraining, which is close in vocabulary but opposite in setting from the
drifting W4A16 inference-time channels measured here.} PMPD changes precision across inference
phases and deeper decode positions \citep{chen2025pmpd}; our current evidence
points in an orthogonal direction, where later decode positions can require
different protected channels rather than simply lower precision.\footnote{We
refer to arXiv:2603.19296 as test-time TTQ when discussing it, to avoid
confusion with classical Trained Ternary Quantization.}

\paragraph{Quantization error accumulation.}
Statistically-Lossless Quantization (SLQ) argues that, in a near-lossless
compression regime, quantization error need not compound across token
positions \citep{helcig2026slq}. That counter-result is important for this
draft because one possible explanation for the failed W4A16 interventions is
long-horizon compound error rather than channel selection alone. The regimes
are not the same: SLQ studies statistically lossless compression, whereas this
draft studies aggressive W4A16 PTQ, small reasoning models, and endpoints up
to 20,000 decode positions. QEP identifies layer-wise error propagation as a
limitation of low-bit PTQ \citep{arai2025qep}, and LAQuant argues that
long-decoding reasoning failures can persist despite short-decode accuracy
unless calibration and lookahead losses account for downstream residual
effects \citep{choi2026laquant}. Activation Sensitivity gives the theoretical
version of the same warning: activation magnitude alone is only a proxy for
loss sensitivity, and cross-layer error accumulation, calibration mismatch,
and task-conditional sensitivity are open challenges for PTQ
\citep{xu2026activationsensitivity}. We therefore treat compound error as an
open hypothesis, not a conclusion. Our KL-accumulation packet directly tests
this mechanism in the aggressive W4A16 endpoint setting and finds flat or
sublinear trajectories rather than superlinear growth, so the paper does not
use compound error as the primary explanation for the observed channel-set
failures.

\paragraph{Channel-wise gated recurrence.}
Recent efficient-recurrence and linear-attention architectures make
data-dependent channel-wise memory control a central design axis. Gated Linear
Attention (GLA) formulates linear attention as recurrent matrix-valued state
with data-dependent gating \citep{yang2023gla}. Gated DeltaNet combines
adaptive memory erasure with delta-rule updates and reports hybrids that mix
Gated DeltaNet with attention or Mamba2 layers \citep{yang2024gdn}. Moonshot
AI's October 2025 Kimi Linear report is a recent example in this line: its
Kimi Delta Attention (KDA) is described as channel-wise gated linear attention
that extends Gated DeltaNet with finer-grained gating, and the released Kimi
Linear model is a layerwise hybrid of KDA and Multi-Head Latent Attention
\citep{kimilinear2025}. The official Qwen3.6 model card likewise describes a
hybrid layout that interleaves Gated DeltaNet blocks with gated attention
\citep{qwen36hf}. These sources motivate OutlierMigrate as an empirical test
of whether high-magnitude channels stay fixed during long decode in hybrid
language models. They do not imply that Kimi Linear, Qwen3.6, GLA, or other
unmeasured models exhibit the migration measured in this draft.

\section{Method}

Phase 0 follows the frozen preregistration at
\url{experimental/outlier_migrate/phase0/preregister_om_phase0.md}. The model
is \texttt{ibm-granite/granite-4.0-h-tiny} at HuggingFace snapshot commit
\texttt{791e0d3d28c86e106c9b6e0b4cecdee0375b6124}. The prompt set contains 12
deterministically selected AIME-2025 traces, indices 0--11, with prompt SHA
\texttt{sha256:2f27c54baa8448e033d6e82f53f775dc6abe38188e4f1e5c0b97e3c74fe7c1dd}
from
\url{experimental/outlier_migrate/phase0/results/om_phase0_20260508T011824Z/prompt_manifest.json}.

Phase 1 follows the frozen preregistration at
\url{experimental/outlier_migrate/phase1/preregister_om_phase1.md}. The model
is \texttt{ibm-granite/granite-4.0-h-small} at HuggingFace snapshot commit
\texttt{b8c0982bab7fde4eb48110f5a069527c008fab39}. The prompt set contains 24
deterministically selected AIME-2025 traces, indices 0--23, with prompt SHA
\texttt{sha256:aa038b29332b6d137d558205ee441163e7ea4cb3cc323eb705a2f5928fd2fe4e}
from
\url{experimental/outlier_migrate/phase1/results/om_phase1_20260508T014959Z/prompt_manifest.json}.

For every transformer layer, the runner records per-channel activation
magnitudes. Phase 0 records decode positions \{100, 500, 1000, 5000, 10000\};
Phase 1 records \{100, 500, 1000, 5000, 10000, 20000\}. Each gate selects the
top 1\% channels by mean magnitude at position 100 and computes the fraction
whose rank at the final preregistered position differs by more than two
channel positions. The reported gate metric is the trace-bootstrap aggregate
over layers with equal layer weighting.

\subsection{Comparison with DecDEC measurement}

DecDEC is the closest prior empirical measurement of decode-time channel-set
change. Its Figure 5 evaluates the recall of static analysis-based outlier
identification against a per-step oracle top-$K$ channel set in short
Transformer decode. Our measurement differs in four ways. First, the horizon
is reasoning-scale: up to 20,000 decode positions rather than roughly 100
decode steps. Second, the model class is reasoning LLMs rather than chat-tuned
Transformers. Third, the architecture set includes hybrid Mamba-2 and
parallel-hybrid models, with DeepSeek-R1-Distill-Qwen-1.5B as a
pure-Transformer reference. Fourth, the prompts are AIME-style long reasoning
traces rather than WikiText/BBH-style evaluation.

The numerical values should not be compared as if they were the same metric.
DecDEC reports low recall for static outlier identification; our main
decomposition reports strict set-leaving from the position-100 top-1\% set to
the final measured top-1\% set, plus within-set rank shuffling. Both
measurements are consistent with a shared warning: static channel identities
can become stale during decode. Our contribution is the long-horizon
reasoning-scale extension and the associated W4A16 method-failure analysis,
not the first observation that outlier-channel sets can change.

\section{Experiments}
\label{sec:experiments}

The Phase 0 result packet is
\url{experimental/outlier_migrate/phase0/results/om_phase0_20260508T011824Z}.
The checker output is
\url{experimental/outlier_migrate/phase0/results/om_phase0_20260508T011824Z/checker_result.json},
the metrics file is
\url{experimental/outlier_migrate/phase0/results/om_phase0_20260508T011824Z/metrics.json},
and artifact completeness is recorded in
\url{experimental/outlier_migrate/phase0/results/om_phase0_20260508T011824Z/artifact_check.json}.
The activation artifact SHA is
\texttt{sha256:2783aa329d82e8f43ea4f342e52f0a25fc283b0b8971880f2da093739474d687}
from
\url{experimental/outlier_migrate/phase0/results/om_phase0_20260508T011824Z/activation_magnitude_manifest.json}.

The Phase 1 result packet is
\url{experimental/outlier_migrate/phase1/results/om_phase1_20260508T014959Z}.
The checker output is
\url{experimental/outlier_migrate/phase1/results/om_phase1_20260508T014959Z/checker_result.json},
the metrics file is
\url{experimental/outlier_migrate/phase1/results/om_phase1_20260508T014959Z/metrics.json},
and artifact completeness is recorded in
\url{experimental/outlier_migrate/phase1/results/om_phase1_20260508T014959Z/artifact_check.json}.
The activation artifact SHA is
\texttt{sha256:326a04351e0dda28cf919fe4745d7d9341011db13d940a06cf9c8470633003e1}
from
\url{experimental/outlier_migrate/phase1/results/om_phase1_20260508T014959Z/activation_magnitude_manifest.json}.

The partial Phase 2 result packet is
\url{experimental/outlier_migrate/phase2/results/om_phase2_nemotron3_20260508T231723Z}.
The checker output is
\url{experimental/outlier_migrate/phase2/results/om_phase2_nemotron3_20260508T231723Z/checker_result.json},
the metrics file is
\url{experimental/outlier_migrate/phase2/results/om_phase2_nemotron3_20260508T231723Z/metrics.json},
and artifact completeness is recorded in
\url{experimental/outlier_migrate/phase2/results/om_phase2_nemotron3_20260508T231723Z/artifact_check.json}.
The model is
\url{nvidia/NVIDIA-Nemotron-3-Nano-30B-A3B-BF16} at snapshot commit
\texttt{cbd3fa9f933d55ef16a84236559f4ee2a0526848}. The activation artifact
SHA is
\texttt{sha256:bc841b938a125f05649f84577c0ddb7ab287213a683892cb81ea6ca06296206d}
from
\url{experimental/outlier_migrate/phase2/results/om_phase2_nemotron3_20260508T231723Z/activation_magnitude_manifest.json}.
This packet is explicitly scoped as partial Nemotron-3-only validation:
Qwen3.6 and Kimi Linear are deferred, and
\texttt{full\_cross\_validation\_complete=false}.

The Phase 3 intervention packet is
\url{experimental/outlier_migrate/phase3/results/om_phase3_20260509T212000Z}.
It follows the frozen preregistration at
\url{experimental/outlier_migrate/phase3/preregister_om_phase3_intervention.md}.
The checker output is
\url{experimental/outlier_migrate/phase3/results/om_phase3_20260509T212000Z/checker_result.json},
the metrics file is
\url{experimental/outlier_migrate/phase3/results/om_phase3_20260509T212000Z/metrics.json},
the per-trace perplexities and recoveries are in
\url{experimental/outlier_migrate/phase3/results/om_phase3_20260509T212000Z/per_trace_metrics.json},
and artifact completeness is recorded in
\url{experimental/outlier_migrate/phase3/results/om_phase3_20260509T212000Z/artifact_check.json}.
The model is again \texttt{ibm-granite/granite-4.0-h-tiny} at snapshot
commit \texttt{791e0d3d28c86e106c9b6e0b4cecdee0375b6124}. The prompt set uses
the fixed 24-trace Phase 1 set with prompt SHA
\texttt{sha256:aa038b29332b6d137d558205ee441163e7ea4cb3cc323eb705a2f5928fd2fe4e}.
The quantization scheme is simple symmetric per-channel INT4 weights with
FP16 activations and BF16 protected-channel rows; no AWQ-style or
SmoothQuant-style activation-aware scaling is used. The scored window is the
64-token window ending at decode position 10000. The preregistered recovery
metric is
\[
1 - \frac{\mathrm{perplexity}_{\mathrm{regime}} -
\mathrm{perplexity}_{\mathrm{BF16}}}{\mathrm{perplexity}_{\mathrm{static1\%}} -
\mathrm{perplexity}_{\mathrm{BF16}}}.
\]
The pass rule is median recovery $\geq 0.50$ with bootstrap 95\% CI lower
bound $>0.30$. The kill rule is recovery $<0.20$ or CI upper $<0.30$.
The mandatory controls are static position-100 top-2\% protection and
magnitude-average top-1\% protection over positions \{100, 1000, 5000,
10000\}.

The preregistered dynamic-outlier pass rule in both gates requires migration
fraction at least 0.05 and bootstrap CI lower bound greater than 0.05. This is
a nonstationarity screen, not a claimed practically meaningful effect-size
standard. Phase 1 uses the same dynamic path because Phase 0 returned
\texttt{PASS\_OM\_PHASE0\_DECODE\_TIME\_MIGRATION}.

\section{Results}

\begin{table}[t]
\centering
\small
\begin{tabular}{>{\raggedright\arraybackslash}p{0.24\linewidth}
                >{\raggedright\arraybackslash}p{0.32\linewidth}
                >{\raggedright\arraybackslash}p{0.32\linewidth}}
\toprule
Quantity & Phase 0 & Phase 1 \\
\midrule
Decision & dynamic-outlier pass & replicated-at-scale pass \\
Strict set-leaving fraction & 0.634244791667 & 0.566234756098 \\
Strict set-leaving 95\% CI & [0.605208333333, 0.664192708333] &
[0.550076219512, 0.581707317073] \\
Within-set rank-shuffling fraction & 0.175260416667 & 0.270934959350 \\
Within-set rank-shuffling 95\% CI & [0.160156250000, 0.191015625000] &
[0.260797764228, 0.280614837398] \\
Original preregistered migration fraction & 0.8178385416666667 & 0.843165650406504 \\
Original migration 95\% CI & [0.797265625, 0.8368489583333334] &
[0.8334349593495936, 0.8511432926829268] \\
Preregistered threshold & $\geq$0.05 and CI lower $>$0.05 &
$\geq$0.05 and CI lower $>$0.05 \\
Artifact complete & true & true \\
Model & \url{ibm-granite/granite-4.0-h-tiny} &
\url{ibm-granite/granite-4.0-h-small} \\
Trace count & 12 & 24 \\
Layer count / hidden size & 40 / 1536 & 40 / 4096 \\
Decode positions & 100, 500, 1000, 5000, 10000 &
100, 500, 1000, 5000, 10000, 20000 \\
Bootstrap seed & 20260508 & 20260508 \\
\bottomrule
\end{tabular}
\caption{OutlierMigrate Phase 0 and Phase 1 gate results. Phase 0 values are
from
\url{experimental/outlier_migrate/phase0/results/om_phase0_20260508T011824Z/checker_result.json}
and
\url{experimental/outlier_migrate/phase0/results/om_phase0_20260508T011824Z/metrics.json};
strict decomposition values are from
\url{experimental/outlier_migrate/phase0/results/om_phase0_20260508T011824Z/migration_decomposition.md}.
Phase 1 values are from
\url{experimental/outlier_migrate/phase1/results/om_phase1_20260508T014959Z/checker_result.json}
and
\url{experimental/outlier_migrate/phase1/results/om_phase1_20260508T014959Z/metrics.json};
strict decomposition values are from
\url{experimental/outlier_migrate/phase1/results/om_phase1_20260508T014959Z/migration_decomposition.md}.}
\label{tab:phase-results}
\end{table}

The original preregistered migration fraction intentionally conflates two
effects: channels can leave the top-1\% set entirely, or they can remain in
the set while shuffling rank by more than two positions. For static channel
protection, the stricter set-membership readout is the more direct
post-hoc interpretability measure, not the preregistered gate criterion.
A channel is counted as strict set-leaving only if it is in the
per-trace, per-layer top-1\% set at decode position 100 and has final rank
greater than the top-1\% boundary at decode position 10000 in Phase 0 or
20000 in Phase 1. Under that definition, 0.634244791667 of Phase 0 base
top-1\% channels and 0.566234756098 of Phase 1 base top-1\% channels leave
the protected set entirely. The within-set shuffling rates are
0.175260416667 and 0.270934959350, respectively. The gate decision remains
based only on the original preregistered migration metric.

The exact Phase 0 provenance values are: model snapshot commit
\url{791e0d3d28c86e106c9b6e0b4cecdee0375b6124}, prompt SHA
\url{sha256:2f27c54baa8448e033d6e82f53f775dc6abe38188e4f1e5c0b97e3c74fe7c1dd},
and activation artifact SHA
\url{sha256:2783aa329d82e8f43ea4f342e52f0a25fc283b0b8971880f2da093739474d687};
these are recorded in
\url{experimental/outlier_migrate/phase0/results/om_phase0_20260508T011824Z/model_provenance.json},
\url{experimental/outlier_migrate/phase0/results/om_phase0_20260508T011824Z/prompt_manifest.json},
and
\url{experimental/outlier_migrate/phase0/results/om_phase0_20260508T011824Z/activation_magnitude_manifest.json}.
The exact Phase 1 provenance values are: model snapshot commit
\url{b8c0982bab7fde4eb48110f5a069527c008fab39}, prompt SHA
\url{sha256:aa038b29332b6d137d558205ee441163e7ea4cb3cc323eb705a2f5928fd2fe4e},
and activation artifact SHA
\url{sha256:326a04351e0dda28cf919fe4745d7d9341011db13d940a06cf9c8470633003e1};
these are recorded in
\url{experimental/outlier_migrate/phase1/results/om_phase1_20260508T014959Z/model_provenance.json},
\url{experimental/outlier_migrate/phase1/results/om_phase1_20260508T014959Z/prompt_manifest.json},
and
\url{experimental/outlier_migrate/phase1/results/om_phase1_20260508T014959Z/activation_magnitude_manifest.json}.

The Phase 0 checker returned
\texttt{PASS\_OM\_PHASE0\_DECODE\_TIME\_MIGRATION}. Its exact artifact reason
was that the point estimate 0.81783854 is at least 0.05 and the CI lower bound
0.79726562 is greater than 0.05. The Phase 1 checker returned
\texttt{PASS\_OM\_PHASE1\_REPLICATED\_AT\_SCALE}. Its exact artifact reason
was that the point estimate 0.84316565 is at least 0.05 and the CI lower bound
0.83343496 is greater than 0.05.

\begin{table}[t]
\centering
\small
\begin{tabular}{>{\raggedright\arraybackslash}p{0.42\linewidth}
                >{\raggedright\arraybackslash}p{0.43\linewidth}}
\toprule
Quantity & Partial Phase 2 Nemotron-3 \\
\midrule
Decision & partial PASS; Qwen3.6/Kimi deferred \\
Strict set-leaving fraction & 0.533713200380 \\
Strict set-leaving 95\% CI & [0.467414529915, 0.599982193732] \\
Within-set rank-shuffling fraction & 0.269082383666 \\
Within-set rank-shuffling 95\% CI & [0.238841405508, 0.299056267806] \\
Original preregistered migration fraction & 0.820809591642925 \\
Original migration 95\% CI & [0.7865325261158594, 0.8544931149097815] \\
Preregistered threshold & $\geq$0.05 and CI lower $>$0.05 \\
Artifact complete & true \\
Model & \url{nvidia/NVIDIA-Nemotron-3-Nano-30B-A3B-BF16} \\
Trace count & 24 \\
Layer count / hidden size & 52 / 2688 \\
Decode positions & 100, 500, 1000, 5000, 10000, 20000 \\
\bottomrule
\end{tabular}
\caption{Partial Phase 2 cross-family validation. The checker-backed packet is
\url{experimental/outlier_migrate/phase2/results/om_phase2_nemotron3_20260508T231723Z}.
This is not full cross-validation because Qwen3.6 and Kimi Linear are deferred.}
\label{tab:phase2-partial}
\end{table}

\begin{table}[t]
\centering
\small
\resizebox{\linewidth}{!}{%
\begin{tabular}{@{}lrrrr@{}}
\toprule
Run & Original drift & Strict set-leaving & Within-set shuffling & Packet \\
\midrule
Granite-Small Phase 1 & 0.836763211382 & 0.566234756098 & 0.270934959350 &
\url{experimental/outlier_migrate/phase9/step9_0_decomposition_replication.md} \\
Nemotron-3 Phase 2 & 0.801667853751 & 0.533713200380 & 0.269082383666 &
\url{experimental/outlier_migrate/phase9/step9_0_decomposition_replication.md} \\
DeepSeek-R1-Distill Phase 5' & 0.834914434524 & 0.670572916667 & 0.165736607143 &
\url{experimental/outlier_migrate/phase9/step9_0_decomposition_replication.md} \\
Falcon-H1 Phase 7 & 0.767887205387 & 0.673611111111 & 0.097537878788 &
\url{experimental/outlier_migrate/phase7/results/om_phase7_falcon_h1_20260512T223600Z/checker_result.json} \\
\bottomrule
\end{tabular}
}
\caption{Cross-architecture decomposition after Phase 5' and Phase 7. Strict
set-leaving remains above 0.50 on Granite-Small, Nemotron-3, DeepSeek, and
Falcon-H1. The Falcon-H1 packet uses post-block residual hooks; it is not
pathway-specific pre-sum attribution.}
\label{tab:cross-arch-decomposition}
\end{table}

Phase 5' broadens the phenomenon beyond measured Mamba-2-style hybrids:
\texttt{deepseek-ai/DeepSeek-R1-Distill-Qwen-1.5B} returns
\texttt{DYNAMIC\_REGIME\_TRANSFORMER} with original drift
0.8393787202380952, 95\% CI [0.8278459821428571, 0.8507254464285714],
strict set-leaving 0.670572916667, and within-set shuffling
0.165736607143. Phase 7 checks Falcon-H1 as a parallel hybrid topology and
returns \texttt{WITHIN\_LINEAGE\_2\_CONSISTENT}: original drift
0.7833543771043772, 95\% CI [0.7544191919191919, 0.8112373737373737], strict
set-leaving 0.673611111111, and within-set shuffling 0.097537878788. These
two additions make the Mamba-specific title untenable and motivate the broader
``decode-position channel drift'' framing.

Experiment D formalizes the decomposition and Experiment E checks threshold
sensitivity. The top-1\% decomposition is stable through top-2\% on all
landed packets. Top-5\% changes the component split and is treated as
sensitivity, not the primary operating point. Phase 9 Step 9.0 additionally
computes adjacent-bin Jaccard overlap for top-channel sets; the overall top-1\%
adjacent-bin Jaccard is 0.745999647119, which made position-binned methods
plausible enough to test.

\subsection{Static and position-conditioned protection attempts}

Phase 3 tests the most direct intervention suggested by the measurement:
instead of protecting only the position-100 top-1\% channels, protect the union
of top-1\% channels observed at decode positions \{100, 1000, 5000, 10000\}.
This is intentionally a static-set intervention. It asks whether membership
migration alone explains enough of the quantization gap to support a simple
deployment recipe.

\begin{table}[t]
\centering
\small
\resizebox{\linewidth}{!}{%
\begin{tabular}{@{}lrrrr@{}}
\toprule
Regime & Mean perplexity & Median perplexity & Median recovery & Recovery 95\% CI \\
\midrule
BF16 baseline & 1.005358821083984 & 1.000282908378305 & reference & reference \\
static-1\% position 100 & 1.015228246226743 & 1.000282342709987 & reference & reference \\
static-2\% position 100 & 1.021467448964073 & 1.000390886047853 & 0.000000000000 & [0.000000000000, 0.356017001423] \\
magnitude average top-1\% & 1.009248547916613 & 1.000236584155715 & 0.061276118929 & [0.000000000000, 0.512245438462] \\
migration-aware union & 1.008573512406127 & 1.000275313560728 & 0.000000000000 & [0.000000000000, 0.711143244199] \\
\bottomrule
\end{tabular}
}
\caption{Phase 3 intervention and controls on Granite-4.0-H-Tiny. Values are
from
\url{experimental/outlier_migrate/phase3/results/om_phase3_20260509T212000Z/per_trace_metrics.json},
\url{experimental/outlier_migrate/phase3/results/om_phase3_20260509T212000Z/metrics.json},
and
\url{experimental/outlier_migrate/phase3/results/om_phase3_20260509T212000Z/control_metrics.json}.
Recovery is defined relative to the BF16-vs-static-1\% gap.}
\label{tab:phase3-intervention}
\end{table}

The intervention kills. The checker returned the decision string
{\small\texttt{KILL\_OM\_PHASE3\_INTERVENTION\_FAILS}} because the union
regime had median recovery 0.000000000000, below the preregistered 0.20 kill
floor. It
also missed the pass rule: median recovery $\geq 0.50$ with CI lower bound
$>0.30$. The artifact packet is complete, and the control-stop condition did
not fire. The best control, magnitude-average protection, exceeded union by
0.061276118929 median recovery, below the human stop threshold of 0.10. Ten
of the 24 traces had no recoverable static gap under this scoring definition.
That fraction, 0.416666666667, also exceeds the preregistered 25\% no-gap kill
condition. The checker reports the first kill reason it reaches, but the
packet satisfies both kill paths.

The recovery statistic is intentionally centered on the per-trace
BF16-vs-static-1\% gap, so it can be negative when a regime is worse than
static-1\%, greater than one when a regime beats BF16 on the scored window, and
zero when there is no recoverable static gap. In the union regime, 11 traces
had positive recovery, 10 were zero because there was no recoverable static
gap, 3 were negative, and 4 exceeded recovery 1.0. This unstable per-trace
surface is part of the negative result rather than a hidden success case: the
primary statistic was preregistered as the median across all 24 traces. The
median was chosen for the gate because the recovery ratio can become very
large when the static gap is small; a mean recovery would be dominated by
unstable high-ratio traces rather than by the typical trace.

\begin{table}[t]
\centering
\small
\resizebox{\linewidth}{!}{%
\begin{tabular}{@{}llrr@{}}
\toprule
Grid & Positions & Median recovery & Recovery 95\% CI \\
\midrule
sparse & \{100, 5000, 10000\} & 0.000000000000 & [0.000000000000, 0.752281234025] \\
primary & \{100, 1000, 5000, 10000\} & 0.000000000000 & [0.000000000000, 0.711143244199] \\
dense & \{100, 500, 1000, 2000, 5000, 7500, 10000\} & 0.000000000000 & [0.000000000000, 0.920007799726] \\
\bottomrule
\end{tabular}
}
\caption{Position-grid sensitivity. Values are from
\url{experimental/outlier_migrate/phase3/results/om_phase3_20260509T212000Z/grid_sensitivity.md}
and
\url{experimental/outlier_migrate/phase3/results/om_phase3_20260509T212000Z/grid_sensitivity_metrics.json}.
Increasing grid density did not produce a nonzero median recovery.}
\label{tab:phase3-grid}
\end{table}

\begin{table}[t]
\centering
\small
\resizebox{\linewidth}{!}{%
\begin{tabular}{@{}lrrrrr@{}}
\toprule
Attempt & Model & Decision & Median recovery & CI95 & Key control readout \\
\midrule
Phase 3 union & Granite-Tiny & \texttt{KILL\_OM\_PHASE3\_INTERVENTION\_FAILS} &
0.000000000000 & [0.000000000000, 0.711143244199] &
10/24 no recoverable static gap \\
Phase 4 union & Granite-Small & \texttt{KILL\_OM\_PHASE4\_INTERVENTION\_FAILS} &
0.000000000000 & [0.000000000000, 0.069540641955] &
9/24 no recoverable static gap \\
Phase 9 M2 bins & Granite-Small & \texttt{KILL\_M2\_RANDOM\_CONTROL\_BEATS} &
-0.866837313391 & [-3.435289251335, 0.595238665766] &
random-bin median beats M2 by 0.667548290168 \\
Phase 9 M10 scales & Granite-Small & \texttt{KILL\_M10\_RANDOM\_CONTROL\_BEATS} &
0.234447927834 & [-2.072905653259, 0.503521494450] &
random-bin median beats M10 by 0.761487459046 \\
Phase 9 M11 EMA & Granite-Small & \texttt{KILL\_M11\_AMBIGUOUS} &
0.048299284138 & [-8.941156151014, 0.704599255873] &
random-walk median is lower by 1.254404832487 \\
Phase 9 M18 act+K & Granite-Small & \texttt{KILL\_M18\_AMBIGUOUS} &
-0.343590844126 & [-14.071191710979, 0.740602124280] &
KIVI key-only median -2.174128311393; random-coupled median -4.664379731051 \\
Phase 9 DecDEC proxy & Granite-Small & \texttt{PASS\_DECDEC\_BASELINE\_REPORTED} &
-0.070034781258 & [-8.495693501233, 0.674176370350] &
random reactive median -1.262604926125; static top-10\% median -0.228273368210 \\
Phase 9 M11b top-5 & Granite-Small & \texttt{PASS\_M11B\_BUDGET\_MATTERS} &
0.449284091125 & [-1.300900018791, 1.000790626547] &
beats static top-10\% by 0.501051267680 median \\
Phase 9 M11b top-5 & Nemotron-3 & \texttt{AMBIGUOUS\_M11B\_NEMOTRON\_BUDGET\_SIGNAL} &
0.456736183270 & [0.345887792388, 0.794986745913] &
does not beat static top-10\%; margin -0.137647560348 \\
Phase 9 M11b top-10 & Nemotron-3 & \texttt{PASS\_M11B\_NEMOTRON\_REPLICATES} &
0.814739798903 & [0.254439288178, 0.922552264880] &
beats static top-10\% by 0.220356055286 median \\
Phase 9 M26 stable core & Granite-Small & \texttt{AMBIGUOUS\_M26} &
0.177615807648 & [-0.116448558090, 0.999915931324] &
beats random matched-core by 1.727084589454 median \\
Phase 9 ParoQuant & Granite-Small & \texttt{PASS\_PAROQUANT\_BASELINE\_REPORTED} &
0.753776848891 & [0.477044452405, 1.003769554323] &
rotation baseline beats M11b top-5 by 0.304492757767 median \\
\bottomrule
\end{tabular}
}
\caption{Intervention attempts. Phase 3 and Phase 4 test static union
protection under W4A16. Phase 9 M2 tests position-conditioned protected sets
A/B/C with a static-3\% matched-cost control and a random-bin negative
control. Phase 9 M10 tests hard-binned scale tables. Phase 9 M11 tests
EMA-smoothed protected-set updates. Phase 9 M18 tests joint activation and
key-cache channel protection. The DecDEC row is a descriptive algorithmic
baseline, not a positive-method gate. M11b is the only channel-set method with
a partial positive signal, but the best budget differs by model and Nemotron
top-5 does not beat matched static top-10. ParoQuant is a rotation baseline,
not a channel-set method. M10 values are from
\url{experimental/outlier_migrate/phase9/results/om_phase9_m10_granite_small_vac12_20260515T085800Z/checker_result.json};
M11 values are from
\url{experimental/outlier_migrate/phase9/results/om_phase9_m11_granite_small_vac12_20260516T010728Z/checker_result.json}; M18 values are from
\url{experimental/outlier_migrate/phase9/results/om_phase9_m18_granite_small_vac12_20260516T193500Z/checker_result.json};
DecDEC values are from
\url{experimental/outlier_migrate/phase9/results/om_phase9_decdec_granite_small_vac12_20260517T141500Z/checker_result.json};
M11b Granite values are from
\url{experimental/outlier_migrate/phase9/results/om_phase9_m11b_granite_small_vac12_reuse_20260518T030300Z/checker_result.json};
M11b Nemotron values are from
\url{experimental/outlier_migrate/phase9/results/om_phase9_m11b_nemotron_static1pct_salvage_20260524T2130Z/checker_result.json};
M26 values are from
\url{experimental/outlier_migrate/phase9/results/om_phase9_m26_granite_small_vac12_20260518T203000Z/checker_result.json};
ParoQuant values are from
\url{experimental/outlier_migrate/phase9/results/om_paroquant_granite_small_20260520T1555Z/checker_result.json}.}
\label{tab:intervention-summary}
\end{table}

Phase 4 repeats the union intervention on Granite-4.0-H-Small with a larger
512-token scoring window. It also kills: median recovery is 0.000000000000,
CI95 [0.000000000000, 0.06954064195470389], and the no-recoverable-static-gap
fraction is 0.375. This rules out the simplest explanation that Phase 3 failed
only because Granite-Tiny was too robust or because the 64-token scoring
window was too short.

Phase 9 M2 then tests a lightly position-conditioned intervention. Instead of
one static union, M2 uses three protected sets: early positions
\{100, 200, 500\}, middle positions \{1000, 2000, 5000\}, and late positions
\{7000, 10000, 15000\}. The active protected set switches at decode positions
800 and 6000. The result kills by the random-control rule. On the
vacation-mode 12-trace deterministic slice, 8 traces have a positive
static-1\% gap. M2 median recovery is -0.866837313391, static-3\% matched-cost
median recovery is -0.625927638959, and random-bin median recovery is
-0.199289023223. The random-bin assignment beats M2 by 0.667548290168 median
recovery, exceeding the preregistered 0.10 kill threshold. The random-bin
control is not good in absolute terms; its importance is that it is less bad
than the intended position-conditioned assignment.

M10 tests whether the problem is not membership but scale choice: it stores
separate SmoothQuant-style scale tables for decode-position bins. This also
kills by the random-control rule. M10 median recovery is 0.234447927834, CI95
[-2.072905653258708, 0.5035214944499808], while the random-bin scale control
has median recovery 0.995935386880 and beats M10 by 0.761487459046. M11 then
removes the hard boundary by replacing bin switches with an EMA-smoothed
protected-set score. M11 does not pass: the best arm, alpha 0.5, has median
recovery 0.048299284138 with CI95 [-8.94115615101423, 0.7045992558730417].
However, its random-walk control has median recovery -1.206105548349, so M11
is structurally different from M2 and M10: smoothing removes the
random-control-beats failure, but smoothing alone does not recover quality.

Together, the intervention results support a three-mechanism account. First,
boundary discontinuities are actively harmful: M2 and M10 are both beaten by
random-bin controls. Second, smoothing removes that harm but does not solve
the gap at top-1\% budget: M11 is not beaten by random-walk, yet misses the
positive-method bar. Third, signal quality and budget separate. M18
activation+K has median recovery -0.343590844126, while KIVI key-only and
random coupled controls have medians -2.174128311393 and -4.664379731051.
That comparison suggests the coupled selection is not arbitrary noise, but it
does not recover quality at top-1\%. DecDEC then tests a reactive
endpoint-selection baseline: median recovery is -0.070034781258 with CI95
[-8.495693501233, 0.674176370350]. It is less damaging than a random reactive
top-1\% control whose median is -1.262604926125, but it does not recover the
W4A16 static gap. M11b changes the picture by increasing budget under a
smoothed update. Granite-Small top-5 recovers 0.449284091125 median and beats
static top-10 by 0.501051267680 median, though the CI is wide. Nemotron
top-5 is also positive with a strictly positive CI, but static top-10 is
stronger; Nemotron top-10 is the arm that beats static top-10. The conclusion
is therefore not "EMA top-5 works" but "budget is load-bearing, and the
budget-selection rule is not yet model-invariant."

A post-M18 diagnostic pass further narrows what the current packets can say.
The K/V-cache migration question is not identifiable from existing artifacts:
only M18 contains K/V-channel evidence, and that evidence is endpoint-only at
decode position 10000 rather than early-to-late K/V snapshots. The recovery
curve question is also not identifiable, because M2, M10, M11, and M18 each
score one 512-token endpoint window. The activation-side evidence is still
useful: Granite-Small fine-grid top-1\% Jaccard changes smoothly across
500-position bins, and all measured models retain a small always-protected
core that is smaller than the full top-1\% set. These diagnostics make DecDEC
and M11b the right next tests rather than another unmotivated method. Those
packets have now landed: DecDEC tests reactive endpoint selection, while M11b
tests whether the remaining failure is simply protection-budget insufficiency.

\subsection{Budget tuning, rotation context, and KL accumulation}

M11b directly tests the budget-insufficiency hypothesis by running the M11 EMA
policy with larger protected-channel budgets. On Granite-Small, the top-5 arm
is the first mechanical channel-set PASS: median recovery is 0.449284091125,
and it beats the static top-10 control by 0.501051267680 median. The
uncertainty is still large, with CI95 [-1.300900018791, 1.000790626547], so
the paper treats it as a signal requiring cross-model validation rather than
as a settled method.

The Nemotron Path C salvage supplies that validation surface after rerunning
the corrected \texttt{static\_1pct} baseline. It resolves the previous
infrastructure failure: \texttt{static\_1pct} now records 6052 quantized tensors and
76 excluded tensors. The result is positive but not a clean top-5 transfer.
Top-5 median recovery is 0.456736183270 with CI95 [0.345887792388,
0.794986745913], but static top-10 recovers 0.594383743618. Top-10 M11b
recovers 0.814739798903 with CI95 [0.254439288178, 0.922552264880] and beats
static top-10 by 0.220356055286 median. This supports the budget mechanism
while showing that the best budget and matched-static comparison are
model-dependent.

ParoQuant provides a different context. Our algorithmic ParoQuant baseline on
Granite-Small recovers 0.753776848891 with CI95 [0.477044452405,
1.003769554323], beating M11b top-5 by 0.304492757767 median. This is not a
contradiction of the channel-drift measurement. It shows that a rotation
baseline can reduce the same W4A16 gap through a different mechanism, so the
paper's negative claims must be scoped to channel-set protection rather than
to W4A16 PTQ generally.

Finally, the KL-accumulation packet weakens the compound-error explanation in
this specific Granite-Small setting. It computes per-position
KL(BF16$\parallel$Q) over 44,208 rows with dense-grid fallback. Mean KL is
0.150302750276 for static-1\%, 0.132701884446 for the DecDEC proxy, and
0.131406585383 for M11. The best fits are flat or sublinear rather than
superlinear, and the fitted AR decay parameters are 0.5147 for static-1\%,
0.4461 for the DecDEC proxy, and 0.4387 for M11. Thus the present evidence
does not support runaway long-horizon compounding as the main failure mode;
signal staleness and budget selection remain the stronger explanations.

\subsection{Layer-stratified migration}

The negative intervention result does not erase the measurement finding: it
narrows the mechanism. Figure~\ref{fig:layer-stratified} uses the existing
Phase 0, Phase 1, and partial Phase 2 activation packets to stratify strict
set-leaving and within-set shuffling by layer type. The result is not confined
to a single Granite mixer family. Granite shows strict set-leaving in both
attention and SSM/Mamba layers, while Nemotron-3 shows it in attention,
SSM/Mamba, and MoE-classified layers. This is consistent with the paper's
diagnostic claim: static membership is unstable in the measured packets, but
successful handling likely needs more than a larger static protected set.

\begin{figure}[t]
\centering
\small
\begin{tabular}{>{\raggedright\arraybackslash}p{0.32\linewidth}
                >{\raggedleft\arraybackslash}p{0.16\linewidth}
                >{\raggedleft\arraybackslash}p{0.17\linewidth}
                >{\raggedleft\arraybackslash}p{0.17\linewidth}}
\toprule
Run / layer type & Layers & Strict set-leaving & Within-set shuffling \\
\midrule
Phase 0 Granite-Tiny attention & 4 & 0.618490 & 0.186198 \\
Phase 0 Granite-Tiny SSM/Mamba & 36 & 0.635995 & 0.174045 \\
Phase 1 Granite-Small attention & 4 & 0.557165 & 0.282266 \\
Phase 1 Granite-Small SSM/Mamba & 36 & 0.567243 & 0.269676 \\
Phase 2 Nemotron-3 attention & 6 & 0.563014 & 0.254630 \\
Phase 2 Nemotron-3 MoE & 23 & 0.526503 & 0.277778 \\
Phase 2 Nemotron-3 SSM/Mamba & 23 & 0.533280 & 0.264157 \\
\bottomrule
\end{tabular}
\caption{Layer-stratified migration decomposition from
\url{experimental/outlier_migrate/phase3/results/layer_stratified_migration.json}.
Values report layer-type means, averaged equally over layers. Strict
set-leaving is the portion most directly relevant to static protected-channel
membership; within-set shuffling motivates scale-refresh mechanisms but is not
itself a membership failure.}
\label{fig:layer-stratified}
\end{figure}

\section{Discussion}

The strongest claim in this draft is now broader and more precise than the
initial Granite-specific hypothesis. Across four measured model families,
channels that are top-1\% by magnitude early in a long decode often do not
remain in the later top-1\% set. The strict set-leaving decomposition is the
most important systems readout because it directly targets the membership
assumption behind static protected-channel maps. The original rank-drift
metric remains the preregistered gate; the stricter decomposition explains why
the gate matters for deployment.

The method results are informative because they separate mechanisms rather
than collapsing every failure into one explanation. Static union protection
addresses only membership; it does not refresh scales, respond to within-set
rank changes, or decide which layers need different treatment. Phase 4 shows
that the Phase 3 failure is not just a Granite-Tiny or 64-token window
artifact. M2 and M10 show that hard protected-set or scale-table binning can
be actively harmful. M11 then separates smoothness from recovery: it removes
the random-control failure but does not restore quality at top-1\%. M18 and
DecDEC show that reactive or cross-tensor signals are better than matched
random controls but still insufficient at top-1\%. M26 shows that the
always-protected stable core has signal but insufficient magnitude. M11b shows
that budget can unlock recovery, while Nemotron shows that the budget optimum
does not simply transfer.

The KL result rules out one tempting explanation in the measured packet. If
compound error were the dominant mechanism, we would expect per-position KL to
grow superlinearly or maintain very high error autocorrelation. Instead, the
trajectories are flat or sublinear and the AR decays are about 0.44--0.51.
This is consistent with the SLQ counter-result in spirit, though in a much
more aggressive W4A16 regime. The paper therefore emphasizes stale or
misbudgeted channel protection rather than runaway accumulation.

ParoQuant constrains the negative claim. A rotation baseline recovers
0.753776848891 median on Granite-Small, which is better than M11b on the same
packet. The paper should therefore be read as evidence that channel-set
protection is fragile in long-decode W4A16 reasoning, not as evidence that
W4A16 PTQ is impossible. The natural ICLR follow-up is not another
threshold-tuned channel set, but composition of rotation-based conditioning
with budget-aware protection, plus broader replication at larger reasoning
model scales.

The theoretical mechanism remains plausible but is not itself an empirical
claim about unmeasured models. In channel-wise gated recurrence or
linear-attention layers, independent per-channel decay rates imply that
channels retain state at different rates. Low-decay channels can accumulate
magnitude over long decode, while high-decay channels shed magnitude. Under
heterogeneous inputs and heterogeneous decay rates, the composition of the top
1\% channels by magnitude should change across decode positions rather than
remain fixed to the position-100 set. Kimi Linear's KDA, Qwen3.6-style Gated
DeltaNet blocks, and GLA-style gated recurrence provide architectural examples
where channel-wise or fine-grained data-dependent gating is central
\citep{kimilinear2025,qwen36hf,yang2024gdn,yang2023gla}. This draft has not
measured Kimi Linear, Qwen3.6, or GLA variants.

\section{Limitations}

This draft is not camera-ready final and should not be read as a deployable
systems method. It reports replicated measurement, failed method classes, and
a partial budget-tuning remedy, not an end-to-end efficiency implementation.
The strongest positive evidence is the cross-model measurement plus the M11b
budget-sensitivity result; the strongest negative evidence is that many
channel-set policies fail or depend strongly on model and budget.

The scope is W4A16 post-training quantization for small reasoning models and
hybrid Mamba-2 or parallel-hybrid architectures, with a small pure-Transformer
reference. It should not be read as an all-formats claim. The underlying
decode-position channel-set drift may be a broader autoregressive phenomenon,
but quality consequences depend on format dynamic range: INT4 exposes severe
failure surfaces here, FP8 may absorb comparable drift in settings like
ThoughtFlow, block formats such as MXFP4/NVFP4 move the analysis from channels
to blocks, and INT2/INT3 would likely be still more fragile. Large standard
reasoning models are also outside the current evidence. Recent reasoning PTQ
work such as ParoQuant reports strong W4A16 recovery on larger reasoning
models \citep{liang2025paroquant}. Our algorithmic ParoQuant Granite-Small
baseline also recovers a large fraction of the BF16-vs-static gap, so the
negative result must stay scoped to channel-set protection methods rather than
to W4A16 PTQ in general.

The empirical surface is still limited. All reported traces are deterministic
AIME-2025 slices. The bootstrap uncertainty is over traces, not over
independent benchmark families, seeds, or prompt distributions. The top-1\%
fraction, rank-delta threshold, and decode-position grid were fixed in the
packets; they should not be retuned on these results. Experiment E shows
reasonable top-1\% to top-2\% stability, but top-5\% changes the component
split.

Cross-lineage validation remains incomplete. Nemotron-3, DeepSeek, and
Falcon-H1 reduce Granite-specific risk, but Qwen3.6 was blocked by hook
accessibility in current SGLang infrastructure without source modification,
and Kimi Linear has not been measured. Falcon-H1 is post-block residual
evidence rather than separate pre-sum SSM/Attention pathway attribution.
Generalization to RWKV-7, GLA-style systems, and other efficient-recurrence
variants is not established.

Finally, the method failures do not prove that decode-position channel drift
is unusable. They prove that static union sets, a larger Granite-Small scoring
window, hard-binned policies, top-1\% EMA smoothing, top-1\% activation/KV
coupling, DecDEC-style reactive selection, and stable-core-only protection are
not enough under this quantization setup. M11b shows that budget tuning can
matter, but Nemotron shifts the strongest arm from top-5 to top-10 and makes
the matched static top-10 control competitive. The next credible evidence
should come from rotation-plus-budget composition, larger-model replication,
and training-time interventions rather than threshold retuning on the same
failed surfaces.

\section{Reproducibility Statement}

The Phase 0 result packet is
\url{experimental/outlier_migrate/phase0/results/om_phase0_20260508T011824Z}.
It includes \texttt{environment.json}, \texttt{model\_provenance.json},
\texttt{prompt\_manifest.json}, \texttt{command\_metadata.json},
\texttt{random\_seed.json}, \texttt{metrics.json},
\texttt{bootstrap\_ci.json}, \texttt{checker\_result.json}, and
\texttt{artifact\_check.json}. The recorded Phase 0 command was:

\begin{verbatim}
experimental/shared/run_phase0_branch.py --branch outlier_migrate
\end{verbatim}

The Phase 1 result packet is
\url{experimental/outlier_migrate/phase1/results/om_phase1_20260508T014959Z}.
It includes the same artifact classes and records the Granite-4.0-H-Small
snapshot commit
\texttt{b8c0982bab7fde4eb48110f5a069527c008fab39}. The recorded Phase 1
command was:

\begin{verbatim}
experimental/outlier_migrate/phase1/run_om_phase1.py
\end{verbatim}

The checker decisions and exact numbers in this draft are taken from the two
Granite result directories above and the partial Phase 2 directory
\url{experimental/outlier_migrate/phase2/results/om_phase2_nemotron3_20260508T231723Z}.
The OutlierMigrate runners emit the artifact-checked metric file
\texttt{metrics.json}; they do not emit a separate
\texttt{profiler\_metrics.json}.
The recorded Phase 2 command was:

\begin{verbatim}
experimental/outlier_migrate/phase2/run_cross_model.py \
  --partial-nemotron3-only \
  --run-id om_phase2_nemotron3_20260508T231723Z \
  --batch-size 4 --dtype bfloat16 --device auto
\end{verbatim}

The Phase 3 intervention packet is
\url{experimental/outlier_migrate/phase3/results/om_phase3_20260509T212000Z}.
It includes \texttt{environment.json}, \texttt{model\_provenance.json},
\texttt{prompt\_manifest.json}, \texttt{command\_metadata.json},
\texttt{random\_seed.json}, \texttt{per\_trace\_metrics.json},
\texttt{metrics.json}, \texttt{control\_metrics.json},
\texttt{grid\_sensitivity\_metrics.json}, \texttt{checker\_result.json}, and
\texttt{artifact\_check.json}. The recorded Phase 3 command was:

\begin{verbatim}
experimental/outlier_migrate/phase3/run_phase3_intervention.py \
  --run-id om_phase3_20260509T212000Z \
  --batch-size 4 --dtype bfloat16 --device auto \
  --reuse-prequant-run-dir \
  experimental/outlier_migrate/phase3/results/om_phase3_20260509T180200Z
\end{verbatim}

The grid-sensitivity artifact was generated with:

\begin{verbatim}
experimental/outlier_migrate/phase3/run_grid_sensitivity.py \
  --run-dir \
  experimental/outlier_migrate/phase3/results/om_phase3_20260509T212000Z
\end{verbatim}

Figure~\ref{fig:layer-stratified} is a LaTeX-native summary table generated
from
\url{experimental/outlier_migrate/phase3/results/layer_stratified_migration.json}.

The Phase 4 intervention packet is
\url{experimental/outlier_migrate/phase4/results/om_phase4_20260511T054000Z}.
The Phase 5' Transformer-control packet is
\url{experimental/outlier_migrate/phase5_prime/results/om_phase5p_20260512T053800Z}.
Experiment D and E artifacts are in
\url{experimental/outlier_migrate/decomposition_analysis/}. The Phase 7
Falcon-H1 packet is
\url{experimental/outlier_migrate/phase7/results/om_phase7_falcon_h1_20260512T223600Z}.
Phase 9 Step 9.0 outputs are
\url{experimental/outlier_migrate/phase9/step9_0_decomposition_replication.md}
and
\url{experimental/outlier_migrate/phase9/step9_0_bin_overlap_analysis.md}.
The Phase 9 M2 packet is
\url{experimental/outlier_migrate/phase9/results/om_phase9_m2_granite_small_vac12_finalized_20260514T233800Z}.
Its checker decision is
\texttt{KILL\_M2\_RANDOM\_CONTROL\_BEATS}; the vacation-mode adaptation notes
are in \url{swarm/vacation_decisions/}.
The Phase 9 M10 packet is
\url{experimental/outlier_migrate/phase9/results/om_phase9_m10_granite_small_vac12_20260515T085800Z};
the Phase 9 M11 packet is
\url{experimental/outlier_migrate/phase9/results/om_phase9_m11_granite_small_vac12_20260516T010728Z};
and the Phase 9 M18 packet is
\url{experimental/outlier_migrate/phase9/results/om_phase9_m18_granite_small_vac12_20260516T193500Z}.
Their checker decisions are
\texttt{KILL\_M10\_RANDOM\_CONTROL\_BEATS},
\texttt{KILL\_M11\_AMBIGUOUS}, and \texttt{KILL\_M18\_AMBIGUOUS}.
Post-M18 analysis reports are in
\url{experimental/outlier_migrate/phase9/post_m18_analysis/}; they document
which mechanism questions are identifiable from the landed packets and which
would require new telemetry.

The DecDEC algorithmic baseline packet is
\url{experimental/outlier_migrate/phase9/results/om_phase9_decdec_granite_small_vac12_20260517T141500Z}.
The M11b Granite packet is
\url{experimental/outlier_migrate/phase9/results/om_phase9_m11b_granite_small_vac12_reuse_20260518T030300Z}.
The M26 stable-core packet is
\url{experimental/outlier_migrate/phase9/results/om_phase9_m26_granite_small_vac12_20260518T203000Z}.
The KL accumulation packet is
\url{experimental/outlier_migrate/phase9/results/om_phase9_kl_granite_small_dense_20260519T085400Z}.
The ParoQuant algorithmic baseline packet is
\url{experimental/outlier_migrate/phase9/results/om_paroquant_granite_small_20260520T1555Z}.
The valid Nemotron M11b Path C salvage packet is
\url{experimental/outlier_migrate/phase9/results/om_phase9_m11b_nemotron_static1pct_salvage_20260524T2130Z}.
The original Nemotron packet is not used for recovery claims because its
\texttt{static\_1pct} baseline was invalid; the Path C packet reruns that baseline
with corrected W4A16 quantization and reuses only validated score caches.

\section{Ethics Statement}

This is an activation-analysis measurement study on benchmark prompts. The
current artifact does not make deployment, user-data, safety, or
systems-efficiency claims. The main ethical risk is overclaiming the replicated
measurement or the partial M11b budget signal as a validated deployable
method. This draft therefore separates measurement, mechanism, baseline
context, and partial-remedy claims explicitly.

\begin{thebibliography}{40}

\bibitem[Arai and Ichikawa(2025)]{arai2025qep}
Yamato Arai and Yuma Ichikawa.
\newblock Quantization Error Propagation: Revisiting Layer-Wise
Post-Training Quantization.
\newblock arXiv:2504.09629, NeurIPS 2025.
\newblock \url{https://arxiv.org/abs/2504.09629}.

\bibitem[Bort(2025)]{bort2025reasoningcosts}
Julie Bort.
\newblock The rise of AI reasoning models is making benchmarking more
expensive.
\newblock TechCrunch, April 10 2025.
\newblock \url{https://techcrunch.com/2025/04/10/the-rise-of-ai-reasoning-models-is-making-benchmarking-more-expensive/}.

\bibitem[Gartner(2026)]{gartner2026inferencecost}
Gartner.
\newblock Gartner predicts that by 2030, performing inference on an LLM with
1 trillion parameters will cost GenAI providers over 90\% less than in 2025.
\newblock Press release, March 25 2026.
\newblock \url{https://www.gartner.com/en/newsroom/press-releases/2026-03-25-gartner-predicts-that-by-2030-performing-inference-on-an-llm-with-1-trillion-parameters-will-cost-genai-providers-over-90-percent-less-than-in-2025}.

\bibitem[MarketsandMarkets(2025a)]{marketsandmarkets2025aiinference}
MarketsandMarkets.
\newblock AI inference market size, share and growth, 2025 to 2030.
\newblock Commercial market report, February 2025.
\newblock \url{https://www.marketsandmarkets.com/Market-Reports/ai-inference-market-189921964.html}.

\bibitem[MarketsandMarkets(2025b)]{marketsandmarkets2025edgeaihardware}
MarketsandMarkets.
\newblock Edge AI hardware industry worth \$58.90 billion by 2030.
\newblock Commercial market report press release, July 4 2025.
\newblock \url{https://www.marketsandmarkets.com/PressReleases/edge-ai-hardware.asp}.

\bibitem[Wang et~al.(2026)]{wang2026mobiquant}
Dongwei Wang, Jinhee Kim, Seokho Han, Denis Gudovskiy, Yohei Nakata,
Tomoyuki Okuno, KhayTze Peong, Kang Eun Jeon, Jong Hwan Ko, Yiran Chen, and
Huanrui Yang.
\newblock MoBiQuant: Mixture-of-Bits Quantization for Token-Adaptive Elastic
LLMs.
\newblock arXiv:2602.20191, February 2026.
\newblock \url{https://arxiv.org/abs/2602.20191}.

\bibitem[Park et~al.(2025)]{park2025decdec}
Yeonhong Park, Jake Hyun, Hojoon Kim, and Jae W. Lee.
\newblock DecDEC: A Systems Approach to Advancing Low-Bit LLM Quantization.
\newblock 19th USENIX Symposium on Operating Systems Design and
Implementation (OSDI 25), 2025.
\newblock \url{https://www.usenix.org/system/files/osdi25-park-yeonhong.pdf}.

\bibitem[Liu et~al.(2024)]{liu2024kivi}
Zirui Liu, Jiayi Yuan, Hongye Jin, Shaochen Zhong, Zhaozhuo Xu, Vladimir
Braverman, Beidi Chen, and Xia Hu.
\newblock KIVI: A Tuning-Free Asymmetric 2bit Quantization for KV Cache.
\newblock ICML 2024.
\newblock \url{https://arxiv.org/abs/2402.02750}.

\bibitem[Liu et~al.(2025)]{liu2025pmkvq}
Tengxuan Liu, Shiyao Li, Jiayi Yang, Tianchen Zhao, Feng Zhou, Xiaohui Song,
Guohao Dai, Shengen Yan, Huazhong Yang, and Yu Wang.
\newblock PM-KVQ: Progressive Mixed-precision KV Cache Quantization for
Long-CoT LLMs.
\newblock arXiv:2505.18610, 2025.
\newblock \url{https://arxiv.org/abs/2505.18610}.

\bibitem[Yang et~al.(2025)]{yang2025attentionpredictor}
Qingyue Yang, Jie Wang, Xing Li, Zhihai Wang, Chen Chen, Lei Chen, Xianzhi
Yu, Wulong Liu, Jianye Hao, Mingxuan Yuan, and Bin Li.
\newblock AttentionPredictor: Temporal Patterns Matter for KV Cache
Compression.
\newblock NeurIPS 2025.
\newblock \url{https://arxiv.org/abs/2502.04077}.

\bibitem[Chen et~al.(2025)]{chen2025pmpd}
Hao Mark Chen, Fuwen Tan, Alexandros Kouris, Royson Lee, Hongxiang Fan, and
Stylianos I. Venieris.
\newblock Progressive Mixed-Precision Decoding for Efficient LLM Inference.
\newblock arXiv:2410.13461, revised February 2025.
\newblock \url{https://arxiv.org/abs/2410.13461}.

\bibitem[Li et~al.(2025)]{li2025qmamba}
Yinglong Li, Xiaoyu Liu, Jiacheng Li, Ruikang Xu, Yinda Chen, and Zhiwei
Xiong.
\newblock QMamba: Post-Training Quantization for Vision State Space Models.
\newblock arXiv:2501.13624, January 2025.
\newblock \url{https://arxiv.org/abs/2501.13624}.

\bibitem[Ramachandran et~al.(2025)]{ramachandran2025ouromamba}
Akshat Ramachandran, Mingyu Lee, Huan Xu, Souvik Kundu, and Tushar Krishna.
\newblock OuroMamba: A Data-Free Quantization Framework for Vision Mamba.
\newblock arXiv:2503.10959, ICCV 2025.
\newblock \url{https://arxiv.org/abs/2503.10959}.

\bibitem[Pierro and Abreu(2024)]{pierro2024mambaptq}
Alessandro Pierro and Steven Abreu.
\newblock Mamba-PTQ: Outlier Channels in Recurrent Large Language Models.
\newblock arXiv:2407.12397, 2024.
\newblock \url{https://arxiv.org/abs/2407.12397}.

\bibitem[Chiang et~al.(2024)]{chiang2024quamba}
Hung-Yueh Chiang, Chi-Chih Chang, Natalia Frumkin, Kai-Chiang Wu, and Diana
Marculescu.
\newblock Quamba: A Post-Training Quantization Recipe for Selective State
Space Models.
\newblock arXiv:2410.13229, 2024.
\newblock \url{https://arxiv.org/abs/2410.13229}.

\bibitem[Chiang et~al.(2025)]{chiang2025quamba2}
Hung-Yueh Chiang, Chi-Chih Chang, Natalia Frumkin, Kai-Chiang Wu, Mohamed S.
Abdelfattah, and Diana Marculescu.
\newblock Quamba2: A Robust and Scalable Post-training Quantization Framework
for Selective State Space Models.
\newblock arXiv:2503.22879, 2025.
\newblock \url{https://arxiv.org/abs/2503.22879}.

\bibitem[Choi et~al.(2026)]{choi2026laquant}
Euntae Choi, Sumin Song, and Sungjoo Yoo.
\newblock LAQuant: A Simple Overhead-free Large Reasoning Model Quantization
by Layer-wise Lookahead Loss.
\newblock arXiv:2605.08755, 2026.
\newblock \url{https://arxiv.org/abs/2605.08755}.

\bibitem[Xu et~al.(2025)]{xu2025mambaquant}
Zukang Xu, Yuxuan Yue, Xing Hu, Zhihang Yuan, Zixu Jiang, Zhixuan Chen,
Jiangyong Yu, Chen Xu, Sifan Zhou, and Dawei Yang.
\newblock MambaQuant: Quantizing the Mamba Family with Variance Aligned
Rotation Methods.
\newblock arXiv:2501.13484, 2025.
\newblock \url{https://arxiv.org/abs/2501.13484}.

\bibitem[Chen and Hemani(2026)]{chen2026quambase}
Yizhi Chen and Ahmed Hemani.
\newblock Late Breaking Results: Quamba-SE: Soft-edge Quantizer for
Activations in State Space Models.
\newblock arXiv:2601.09451, DATE Late Breaking Results 2026.
\newblock \url{https://arxiv.org/abs/2601.09451}.

\bibitem[Xiao et~al.(2023)]{xiao2023smoothquant}
Guangxuan Xiao, Ji Lin, Mickael Seznec, Hao Wu, Julien Demouth, and Song Han.
\newblock SmoothQuant: Accurate and Efficient Post-Training Quantization for
Large Language Models.
\newblock arXiv:2211.10438, ICML 2023.
\newblock \url{https://arxiv.org/abs/2211.10438}.

\bibitem[Lin et~al.(2023)]{lin2023awq}
Ji Lin, Jiaming Tang, Haotian Tang, Shang Yang, Wei-Ming Chen, Wei-Chen Wang,
Guangxuan Xiao, Xingyu Dang, Chuang Gan, and Song Han.
\newblock AWQ: Activation-aware Weight Quantization for LLM Compression and
Acceleration.
\newblock arXiv:2306.00978, MLSys 2024.
\newblock \url{https://arxiv.org/abs/2306.00978}.

\bibitem[Ashkboos et~al.(2024)]{ashkboos2024quarot}
Saleh Ashkboos, Amirkeivan Mohtashami, Maximilian L. Croci, Bo Li, Pashmina
Cameron, Martin Jaggi, Dan Alistarh, Torsten Hoefler, and James Hensman.
\newblock QuaRot: Outlier-Free 4-Bit Inference in Rotated LLMs.
\newblock arXiv:2404.00456, 2024.
\newblock \url{https://arxiv.org/abs/2404.00456}.

\bibitem[Hooper et~al.(2024)]{hooper2024kvquant}
Coleman Hooper, Sehoon Kim, Hiva Mohammadzadeh, Michael W. Mahoney, Yakun
Sophia Shao, Kurt Keutzer, and Amir Gholami.
\newblock KVQuant: Towards 10 Million Context Length LLM Inference with KV
Cache Quantization.
\newblock arXiv:2401.18079, NeurIPS 2024.
\newblock \url{https://arxiv.org/abs/2401.18079}.

\bibitem[Jang and Tambe(2025)]{jang2025blockdialect}
Wonsuk Jang and Thierry Tambe.
\newblock BlockDialect: Block-wise Fine-grained Mixed Format Quantization for
Energy-Efficient LLM Inference.
\newblock arXiv:2501.01144, ICML 2025.
\newblock \url{https://arxiv.org/abs/2501.01144}.

\bibitem[Helcig et~al.(2026)]{helcig2026slq}
Michael Helcig, Eldar Kurtic, and Dan Alistarh.
\newblock Statistically-Lossless Quantization of Large Language Models.
\newblock arXiv:2605.02404, 2026.
\newblock \url{https://arxiv.org/abs/2605.02404}.

\bibitem[Liang et~al.(2025)]{liang2025paroquant}
Yesheng Liang, Haisheng Chen, Song Han, and Zhijian Liu.
\newblock ParoQuant: Pairwise Rotation Quantization for Efficient Reasoning
LLM Inference.
\newblock arXiv:2511.10645, 2025.
\newblock \url{https://arxiv.org/abs/2511.10645}.

\bibitem[Liao and Wen(2026)]{liao2026chanmix}
Chengxi Liao and Zeyi Wen.
\newblock Channel-Aware Mixed-Precision Quantization for Efficient
Long-Context Inference.
\newblock ICLR 2026.
\newblock \url{https://openreview.net/forum?id=yjr2jX41qO}.

\bibitem[Qwen Team(2026)]{qwen36hf}
Qwen Team.
\newblock Qwen3.6-35B-A3B model card.
\newblock Hugging Face, 2026.
\newblock \url{https://huggingface.co/Qwen/Qwen3.6-35B-A3B}.

\bibitem[Xu(2026)]{xu2026activationsensitivity}
Bruce Changlong Xu.
\newblock Activation Sensitivity as a Unifying Principle for Post-Training
Quantization.
\newblock arXiv:2601.11663, 2026.
\newblock \url{https://arxiv.org/abs/2601.11663}.

\bibitem[Yang et~al.(2023)]{yang2023gla}
Songlin Yang, Bailin Wang, Yikang Shen, Rameswar Panda, and Yoon Kim.
\newblock Gated Linear Attention Transformers with Hardware-Efficient
Training.
\newblock arXiv:2312.06635, 2023.
\newblock \url{https://arxiv.org/abs/2312.06635}.

\bibitem[Yang et~al.(2024)]{yang2024gdn}
Songlin Yang, Jan Kautz, and Ali Hatamizadeh.
\newblock Gated Delta Networks: Improving Mamba2 with Delta Rule.
\newblock arXiv:2412.06464, 2024.
\newblock \url{https://arxiv.org/abs/2412.06464}.

\bibitem[Moonshot AI Kimi Team(2025)]{kimilinear2025}
Moonshot AI Kimi Team: Yu Zhang et~al.
\newblock Kimi Linear: An Expressive, Efficient Attention Architecture.
\newblock arXiv:2510.26692, October 2025.
\newblock \url{https://arxiv.org/abs/2510.26692}.

\bibitem[Zhang et~al.(2025)]{zhang2025mixkvq}
Tao Zhang, Ziqian Zeng, Hao Peng, Huiping Zhuang, and Cen Chen.
\newblock MixKVQ: Query-Aware Mixed-Precision KV Cache Quantization for
Long-Context Reasoning.
\newblock arXiv:2512.19206, 2025.
\newblock \url{https://arxiv.org/abs/2512.19206}.

\end{thebibliography}

\end{document}
